\documentclass[11pt,a4paper]{article}

\usepackage[T2A]{fontenc}
\usepackage[utf8]{inputenc}
\usepackage[russian]{babel}
\usepackage{amsmath,amssymb,bm}
\usepackage[a4paper,margin=2cm]{geometry}
\usepackage{microtype}

\newcommand{\R}{\mathbb{R}}
\newcommand{\T}{\mathsf{T}}
\newcommand{\Id}{I}

\title{Задача оптимизации в Average Derivative Procedure}
\author{}
\date{}

\begin{document}
\maketitle
\vspace{-1.5em}

Рассматривается одноиндексная модель
\[
Y_i=f(\bm\beta^{\T}X_i)+\varepsilon_i,
\qquad
X_i\in\R^d,
\qquad
\|\bm\beta\|_2=1.
\]
Для локального центра $x_j$ задаются веса $w_{ij}$ и их взвешенное среднее
\[
M_j=\frac{\sum_i w_{ij}X_i}{\sum_i w_{ij}}.
\]
Центрирование относительно $M_j$ даёт
\[
\sum_i w_{ij}(X_i-M_j)=0.
\]
Поэтому свободный член локальной линейной аппроксимации не входит в направленные статистики: в задаче ADP можно положить $c_j=0$.

Для набора направлений $\{\varphi_{jr}\}_{r=1}^{P}$ определим
\[
I_{jr}
=
\sum_i Y_i\,\langle X_i-M_j,\varphi_{jr}\rangle w_{ij},
\]
\[
U_{jr}
=
\sum_i (X_i-M_j)
\langle X_i-M_j,\varphi_{jr}\rangle w_{ij}
\in\R^d.
\]
Положим
\[
\bm I_j=(I_{j1},\ldots,I_{jP})^{\T}\in\R^P,
\qquad
U_j=
\begin{pmatrix}
U_{j1}^{\T}\\
\vdots\\
U_{jP}^{\T}
\end{pmatrix}
\in\R^{P\times d}.
\]
Локальная линеаризация функции связи приводит к модели
\[
\bm I_j\approx \ell_jU_j\bm\beta,
\]
где $\ell_j$ аппроксимирует локальную производную
$f'(\bm\beta^{\T}M_j)$.

\section*{Совместная задача}
При фиксированных весах и статистиках внутренняя задача ADP имеет вид
\begin{equation}
\min_{\substack{\bm\beta\in\R^d,\;\|\bm\beta\|_2=1\\
\ell_1,\ldots,\ell_J\in\R}}
\frac12\sum_{j=1}^{J}
\|\bm I_j-\ell_jU_j\bm\beta\|_2^2.
\label{eq:joint}
\end{equation}
Произведение $\ell_jU_j\bm\beta$ билинейно по паре $(\ell_j,\bm\beta)$, поэтому задача \eqref{eq:joint} совместно невыпукла. При фиксации одного блока она становится выпуклой квадратичной задачей по другому блоку. Тем самым \eqref{eq:joint} имеет бивыпуклую структуру.

Кроме того, модель неидентифицируема по масштабу:
\[
\bm\beta\mapsto a\bm\beta,
\qquad
\ell_j\mapsto a^{-1}\ell_j.
\]
Ограничение $\|\bm\beta\|_2=1$ устраняет эту неоднозначность с точностью до знака $\bm\beta\sim-\bm\beta$.

\section*{Альтернированные подзадачи}
При фиксированной $\bm\beta$ каждая локальная задача независима:
\[
\widehat\ell_j
=
\arg\min_{\ell\in\R}
\|\bm I_j-\ell U_j\bm\beta\|_2^2.
\]
Её замкнутое решение равно
\begin{equation}
\widehat\ell_j
=
\frac{\langle\bm I_j,U_j\bm\beta\rangle}
{\|U_j\bm\beta\|_2^2}.
\label{eq:l-update}
\end{equation}
При малом знаменателе применяется ridge-модификация
\[
\widehat\ell_j
=
\frac{\langle\bm I_j,U_j\bm\beta\rangle}
{\|U_j\bm\beta\|_2^2+\eta},
\qquad \eta>0.
\]

При фиксированных $(\ell_j)_{j=1}^{J}$ рассматривается регуляризованная подзадача
\begin{equation}
\min_{\bm\beta\in\R^d}
\frac12\sum_{j=1}^{J}
\|\bm I_j-\ell_jU_j\bm\beta\|_2^2
+
\frac{\lambda}{2}
\|\bm\beta-\bm\beta_{\mathrm{init}}\|_2^2.
\label{eq:beta-subproblem}
\end{equation}
Её гессиан
\[
H=\sum_{j=1}^{J}\ell_j^2U_j^{\T}U_j+\lambda\Id_d
\]
положительно полуопределён при $\lambda=0$ и положительно определён при $\lambda>0$. Условие первого порядка имеет вид
\[
H\widehat{\bm\beta}=b,
\qquad
b=
\sum_{j=1}^{J}\ell_jU_j^{\T}\bm I_j
+
\lambda\bm\beta_{\mathrm{init}}.
\]
Следовательно,
\begin{equation}
\widehat{\bm\beta}
=
\left(
\sum_{j=1}^{J}\ell_j^2U_j^{\T}U_j+\lambda\Id_d
\right)^{-1}
\left(
\sum_{j=1}^{J}\ell_jU_j^{\T}\bm I_j
+\lambda\bm\beta_{\mathrm{init}}
\right).
\label{eq:beta-update}
\end{equation}
Формула \eqref{eq:beta-update} представляет нормальные уравнения ridge-регрессии. В реализации обратная матрица не строится: решается система $H\widehat{\bm\beta}=b$.

После свободного решения выполняется нормировка
\[
s=\|\widehat{\bm\beta}\|_2,
\qquad
\bm\beta\leftarrow\widehat{\bm\beta}/s,
\qquad
\ell_j\leftarrow s\ell_j,
\]
которая сохраняет произведения $\ell_jU_j\bm\beta$. Такая проекция не совпадает с точным решением квадратичной задачи непосредственно на сфере. Для точного учёта ограничения $\|\bm\beta\|_2=1$ требуется сферический, или риманов, решатель.

\end{document}
